% GDZ PPN236018647,native571middle–572upper;HKUST572–573.
\NumberedItem{98}
\indent Terminum medium arithmetico-geometricum inter \(1\) et \(\surd2\) esse \(=\frac\pi\varpi\) usque ad figuram undecimam comprobavimus, qua re demonstrata prorsus novus campus in analysi certo aperietur.
\par\noindent\hfill Br[unsvigae, 1799] Mai. 30.

\medskip\begingroup\small
\makeatletter\tagsleft@true\let\veqno\@@leqno\makeatother
\renewcommand{\thefootnote}{**)}
\indent Der hier niedergelegten Bemerkung, sie mag dem Zufall oder einer Vermutung ihre Entstehung verdanken, muß die genaue Berechnung der beiden Größen \(\frac\pi\varpi\) und \(M(\sqrt2,1)\) vorhergegangen sein. Für \(\frac\pi\varpi\) läßt sich dies auch in der Tat nachweisen. Nach der \textsc{Euler}schen Beziehung\footnote{Siehe \textit{Institutiones calculi integralis} I (1768), § 334, Opera omnia, ser. I, vol. 11, S. 210, vergl. oben S. 150, Gl. (3), wo infolge eines Druckfehlers \(\frac\pi2\) statt \(\frac\pi4\) steht.}
\[
\int_0^1\frac{dx}{\sqrt{1-x^4}}\cdot\int_0^1\frac{x^2\,dx}{\sqrt{1-x^4}}=\frac\pi4
\]
\noindent st nämlich
\begin{equation*}
\tag*{(*)}2\int_0^1\frac{x^2\,dx}{\sqrt{1-x^4}}=\frac\pi\varpi,
\end{equation*}
\noindent da ja
\[
\varpi=2\int_0^1\frac{dx}{\sqrt{1-x^4}}
\]
\renewcommand{\thefootnote}{***)}
\noindent ist. Den bei \textsc{Stirling}\footnote{\textit{Methodus Differentialis etc.}, Londini 1730, S. 57.} auf 17 Dezimalstellen berechneten Wert des Integrals \(\displaystyle\int_0^1\frac{x^2\,dx}{\sqrt{1-x^4}}\), d. h. der \textit{Ordinata Curvae Elasticae} hat \textsc{Gauss} damals sicherlich gekannt (vergl. oben S. 145); auf das Schutzblatt des \textsc{Leiste} hat er neben den dort notierten Näherungswert \(1{,}198\pm\) des Integrales (*) mit genau denselben Schriftzügen, wie sie die Tagebuchnotiz Nr. 98 zeigt, die Zeichen »\(=\frac\pi\varpi\)« hingeschrieben (siehe oben S. 146,
\par\endgroup

\clearpage\SourceLeaf{572}\ArithmeticHead{543}{R}{1799 APR. 8—1799 SEPT.}
\begingroup\small
\noindent art. [3.]). Ferner tritt die Größe \(\frac\pi\varpi\) auf, wenn man aus der Gleichung
\[
\sin\text{ lemn }\varphi=\sin\text{ circuli }\chi
\]
\noindent (siehe oben S. 170, art. [6.]) \(\chi\) durch \(\varphi\) in Bogenmaß darstellen will. Endlich haben die Darstellungen der Zähler und Nenner von sinus lemniscaticus, \(P\varphi,Q\varphi\), wie sie in der Scheda Aa (Werke III, S. 418 und oben S. 537) und auch in der Scheda Ac (siehe oben S. 195, Gleichung [8], [9]) gegeben sind, beide den Faktor \(\sqrt{\frac\pi\varpi}\). — In bezug auf \(M(\sqrt2,1)\) wissen wir nur, daß es als Beispiel in der Scheda Ab (siehe oben S. 174) und auch im \textsc{Leiste} (siehe oben S. 180, dort allerdings \(M(1,\sqrt{\frac12})\)) auftritt, aber diese beiden Stellen sind wohl erst nach der in der Tagebuchnotiz Nr. 98 gemachten Bemerkung geschrieben. Die Veranlassung \(M(\sqrt2,1)\) zu betrachten, könnte für \textsc{Gauss} vielleicht in einer Bemerkung \textsc{Stirlings} gelegen haben, der auf S. 57 seines angeführten Werkes im Exemplum IV sagt: »Quod si longitudini Elasticae adiiciatur sua Ordinata, habebitur numerus \(1{,}\ 9100\ 9889\ 4513\ 8559\ 8\), qui est semiperipheria Ellipseos habentis \(1\) et \(\sqrt2\) pro Axibus«. Es ist in der Tat
\[
\int_0^1\frac{dx}{\sqrt{1-x^4}}+\int_0^1\frac{x^2\,dx}{\sqrt{1-x^4}}=\int_0^1 dx\sqrt{\frac{1+x^2}{1-x^2}}\mathpunct{\raisebox{.35ex}{.}}
\]
\noindent In der Abhandlung von 1800 hat \textsc{Gauss} im Exemplum 4. den Wert von \(M(\sqrt2,1)\) auf 19 Dezimalstellen berechnet (siehe Werke III, S. 364).

\indent Die Voraussage, daß der Beweis für dieses Ineinandergreifen der lemniskatischen Funktionen mit dem agM. sicher ein neues Feld der Analysis eröffnen werde, zeigt, eine wie hohe Bedeutung \textsc{Gauss} der Bemerkung vom 30. Mai 1799 beigelegt hat. Dies kommt auch darin zum Ausdruck, daß er diese Bemerkung an \textsc{Joh. Friedr. Pfaff} brieflich mitgeteilt hat. Aus der oben S. 232 abgedruckten Antwort \textsc{Pfaffs} vom 24. November 1799 geht hervor, daß \textsc{Gauss}, als er an \textsc{Pfaff} über den Gegenstand schrieb, noch nicht im Besitz des Beweises der Gleichung
\[
\frac\pi\varpi=M(\sqrt2,1)
\]
\noindent gewesen ist. Die vollständige Durchführung dieses Beweises scheint ihm erst im November 1799 gelungen zu sein (siehe die Bemerkungen zu den Abschnitten [II.] bis [IV.] der \textit{Theorie des agM.}, oben S. 273 und die Tagebuchnotiz Nr. 100). Die Tragweite der Bemerkung der Nr. 98 für die Entwicklung von \textsc{Gauss}’ Gedankengang besteht darin, daß \textsc{Gauss} durch sie auf die Wichtigkeit des \textls[120]{reziproken Wertes} des agM. aufmerksam wurde, während er früher (vergl. das \textit{Specimen}, oben S. 172) nur das agM. selbst betrachtet hatte. Vergl. dafür auch den Schluß des art. 5. der Abhandlung von 1800, Werke III, S. 366.
\par\noindent\hfill\textsc{Klein.\quad Schlesinger.}
\par\endgroup
